\chapter{Evaluation}

In the project proposal, four success criteria were laid out.




The first three were evaluated as part of a usability study, described in section 4.1. Criterion 4 was evaluated using some timing analysis, which is discussed in section 4.2.

\section{Study}

This section details the study performed to evaluate the app. Section 4.1.1 does....

\subsection{Research Questions}

With the set of success criteria in mind, the following research questions were established to evaluate the overall usability and efficacy of the app.

\subsubsection{Efficacy}
\textbf{RQ 1: }To what extent does the tool help users to learn coding? 

This RQ aims to evaluate the tool's efficacy, in particular by assessing whether users perceive that the tool helps them learn.

\textbf{RQ 2: }Are the challenges provided by the tool appropriately scoped for beginner learners?

This RQ aims to directly address criterion 3 by assessing whether the tool is appropriately scoped to beginners.

\subsubsection{Usability}
\textbf{RQ 3: }How usable is the tool for beginners?

This RQ aims to address criterion 1 by assessing whether beginners were able to understand how to use the tool without help.

\textbf{RQ 4: }How usable is the tool as a simple code editor?

This RQ aims to address criterion 2 by assessing whether the tool can transition with users as they learn, eventually leaving them competent with a standard code editor.

\subsection{Study Design}

The study described below was approved by the Department of Computer Science and Technology Ethics Committee.

The questionnaire consisted of two main sections. The chosen questions needed to accurately measure what they claimed. For the usability section, the well-established System Usability Scale \cite{brooke_sus_1995} was used. It was suitable because...

\subsection{Results}


\subsection{Discussion}

\section{Timing stuff}